% !TEX program = xelatex
\documentclass[11pt,a4paper]{article}

% ---------- Packages ----------
\usepackage[a4paper,margin=1.65cm]{geometry}
\usepackage{titlesec}
\usepackage{enumitem}
\usepackage[hidelinks]{hyperref}
\usepackage{fontawesome5}
\usepackage{fontspec}
\usepackage{graphicx}
\usepackage{xcolor}
\setmainfont{Times New Roman}

% ---------- Style ----------
\pagenumbering{gobble}
\setlength{\parindent}{0pt}
\setlength{\parskip}{2pt}
\setlist[itemize]{leftmargin=1.2em,itemsep=0.5pt,topsep=1pt,parsep=0pt,partopsep=0pt}
\definecolor{Accent}{HTML}{0E3A5D}
\definecolor{SoftGray}{HTML}{666666}

\titleformat{\section}
  {\large\bfseries\color{Accent}}{}
  {0pt}{}
  [\vspace{-4pt}\titlerule\vspace{2pt}]
\titlespacing*{\section}{0pt}{7pt}{3pt}

\newcommand{\cvitem}[2]{\textbf{#1}: #2}
\newcommand{\contactsep}{\hspace{0.35em}{\color{SoftGray}\textbar}\hspace{0.35em}}

\newcommand{\cventry}[4]{%
  \textbf{#1}\hfill \textbf{#2}\\
  \textit{#3}\hfill \textit{#4}\\[-4pt]
}

\newcommand{\cvproject}[4]{%
  \textbf{#1}\hfill \textbf{#2}\\
  \textit{#3}\hfill \textit{#4}\\[-6pt]
}

\begin{document}

% ---------- Header ----------
\begin{minipage}[c]{0.74\textwidth}
{\fontsize{24}{26}\selectfont\textbf{WANG MINGYANG}}\\[-1pt]
\small

{\color{SoftGray}Tokyo, Japan}
\contactsep \faPhone\ \href{tel:07065219182}{070-6521-9182}
\contactsep \faEnvelope\ \href{mailto:wangmingyang_career@outlook.com}{wangmingyang\_career@outlook.com}\\[-1pt]
\faGithub\ \href{https://github.com/AWANGMY}{github.com/AWANGMY}
\contactsep \textbf{Research Field:} AI / Sport Science / Biomechanics
\end{minipage}
\hfill
\begin{minipage}[c]{0.23\textwidth}
\raggedleft
\IfFileExists{WANG(black).jpg}{\includegraphics[width=2.2cm]{WANG(black).jpg}}{}
\end{minipage}


% ---------- Education ----------
\section*{EDUCATION}

{\setlength{\parskip}{1pt}
\cventry
{Dalian University of Technology}{Dalian, China}
{B.Eng in Mechanical Design, Manufacturing and Automation (Japanese Intensive Class)}{09.2020 -- 06.2025}
{\small \textbf{GPA:} 3.55/4.0 \contactsep \textbf{Award:} Excellent Graduation Thesis}\\[-2pt]
\vspace{-3pt}

\cventry
{Tokyo Institute of Technology}{Tokyo, Japan}
{Young Scientists Exchange Program (YSEP), Department of Mechanical Engineering}{10.2023 -- 08.2024}
\vspace{-3pt}

\cventry
{University of Tokyo}{Tokyo, Japan}
{Master's program, Graduate School of Engineering, Mechanical Engineering}{10.2025 -- Present}
}

% ---------- Experience ----------
\section*{EXPERIENCE}

\cvproject
{A Wearable Sensor Based Musculoskeletal Digital Twin System}{Dalian, China}
{Dalian University of Technology}{09.2024 -- 06.2025}
\vspace{-5pt}
\begin{itemize}
  \item Proposed a musculoskeletal digital twin modeling method integrating wearable IMU sensors with the OpenSim platform for pre-TKA gait evaluation.
  \item Collected gait data with Vicon optical motion capture and wearable IMU, and obtained ground reaction force data from a Bertec instrumented treadmill.
  \item Reconstructed 3D knee joint angles, joint moments, and major muscle forces across three anatomical planes via model scaling, inverse kinematics, and inverse dynamics.
  \item Comparative experiments showed high consistency of IMU with optical motion capture, with stronger flexibility in occluded or space-constrained environments.
\end{itemize}

\cvproject
{Skeletal Muscle Contraction Control Utilizing Inverse Models Considering Muscle Fatigue}{Tokyo, Japan}
{Tokyo Institute of Technology}{10.2023 -- 08.2024}
\vspace{-5pt}
\begin{itemize}
  \item Simulated the stimulation voltage required to generate target muscle force using a muscle contraction inverse model.
  \item Incorporated muscle fatigue into simulation calculations to estimate stimulation voltage required to sustain muscle force over long durations.
  \item Conducted electrical stimulation experiments on frog thigh muscle to validate the simulation results.
\end{itemize}

\cvproject
{Human Fatigue Prediction with Multi RGB-D Motion Analysis}{Tokyo, Japan}
{University of Tokyo}{10.2025 -- Present}

\vspace{-5pt}
\begin{itemize}
  \item Built a multi-camera RGB-D pipeline for human keypoint extraction, synchronization, and 3D trajectory reconstruction.
  \item Implemented motion filtering modules, including adaptive Kalman filtering and temporal denoising, to improve signal stability.
  \item Extracting movement features and training machine-learning models to predict fatigue level during human motion tasks.
  \item Developing reproducible workflows for dataset management, model evaluation, and result visualization.
\end{itemize}

% ---------- Skills ----------
\section*{SKILLS}

\cvitem{Programming}{Python, MATLAB, C}\\
\cvitem{Computer Vision \& ML}{Multi-camera calibration, Kalman filtering, CNN}\\
\cvitem{Engineering Tools}{OpenSim, Inventor}\\
\cvitem{Languages}{Chinese (Native), English (TOEFL 97), Japanese (JLPT N1)}\\
\cvitem{Activities}{University Imaging Studio member; Career Planning and Employment Association member; Overseas Exchange Student Ambassador}\\
\cvitem{Interests}{Soccer, basketball, table tennis, photography, cooking}

\end{document}
